\documentclass{article}
\usepackage[utf8]{inputenc}
\usepackage[english]{babel}
\usepackage[margin=1in]{geometry}
\usepackage[]{amsthm} 
\usepackage[]{amsmath} 
\usepackage[]{amssymb} 

\title{Multicast file sharing -- Report}
\author{Kevin Wu, Sam Zhang}
\date{Due: April 6, 2026}

\begin{document}
\maketitle

\setlength{\parindent}{0ex} 
\setlength{\parskip}{5pt} % Adjusts the space between paragraphs

\section*{Overview}


 We introduce a Negative Acknowledgement (NACK)-oriented reliable multicast file-sharing protocol that supports full file transfer for both existing receivers and late joiners with low network and host overhead. The sender is effectively stateless and cyclically transmits control packets containing metadata for all files. Receivers are stateful: they track received chunks and request missing chunk ranges using NACK packets. Integrity is verified at two levels: a per-packet chunk checksum during reception and a file-level checksum upon completion. This protocol provides high reliability and fast catch-up for new joiners, but in worst-case receiver-swarm scenarios it can generate substantial repair traffic.

\section*{Protocol Design}

% Protocl Description, metadata packet, data  packet, NACK packet


\subsection*{Sender}

\subsection*{Receiver}

\subsection*{Comparison to FLUTE}

A different multicast protocl is File Delivery over Unidirectional Transport (FLUTE), where the sender does not have a receive channel. Instead, it cyclically sends all data chunks perpetually. Reliability is ensured by receivers' eventual consistency, as they will receive any missing chunks in the next send cycle.

We preferred the simplicity of FLUTE but did not pick this protocol as the network load is although less volatile, is higher especially on situations where minimal amount of new data transmission is needed, such as when there are no current new joiners, or some receivers are only missing couple chunks. As one of the requirement is to minimize network load by minimizing control packets and data packets. Therefore, we use our current protocol.

\section*{Integrity Verification}

\subsection*{Chunk-Level Checksum}

Chunk-level checksum is computed at data packet generation time on the sender side and carried with each DATA packet. On the receiver side, every incoming chunk is validated against that checksum before being accepted. If the chunk-level checksum fails, we treat it as a bad chunk and immediately re-request that chunk range using a NACK.

\subsection*{File-Level Combined Checksum}

File-level checksum is computed on sender initialization (before transmission starts) and included in metadata. After a receiver has all chunks for a file, it computes the combined file checksum locally and compares it with the expected value from metadata.

This gives us a two-step guard to ensure files are intact: first at chunk granularity during reception, then at whole-file granularity at completion. If file-level verification fails, we work under the assumption that individual chunks were fine when received and that some local corruption happened after/during assembly. Since it is difficult to reliably pinpoint which local chunk(s) got corrupted in that case, we NACK the whole file range and re-fetch the full file.




\section*{Performance Benchmark}

\section*{Limitations}

\end{document}
